\title{Group Sequence Policy Optimization}

\begin{abstract}
We present Group Sequence Policy Optimization (GSPO), a robust and efficient reinforcement learning method tailored for training large language models. Distinct from prior techniques that rely on token-level importance ratios, GSPO employs sequence-level likelihoods to define importance ratios and conducts all reward computations, clipping, and optimization at the sequence level. Our results indicate that GSPO offers significant advantages over the GRPO algorithm, including higher training efficiency and improved performance, alongside enhanced stability for Mixture-of-Experts (MoE) RL training. Additionally, GSPO has the potential to streamline the development of RL infrastructure. These capabilities have substantially contributed to the advancements found in the most recent Qwen3 models.
\end{abstract}

\section{Introduction}
\label{sec:intro}

Reinforcement learning (RL) has become a crucial approach for advancing the capabilities of language models at scale \citep{o1,dpsk-r1,qwq32b,qwen3}. By leveraging large-scale RL, these models acquire the ability to address complex tasks—including competition-level mathematics and programming—by engaging in more intricate and extended forms of reasoning.

Achieving efficient scaling of RL through increased computational resources necessitates, above all, the preservation of stable and reliable training dynamics. Nevertheless, contemporary leading RL methods, such as GRPO \citep{grpo}, suffer from pronounced instability when applied to extremely large language models, frequently leading to severe, irreversible collapse of model performance \citep{qwen3, minimax-m1}. Such training instabilities pose a significant barrier to advancing language model capabilities through further RL-based refinement.

This study demonstrates that the root cause of GRPO's instability is the improper use and subsequent breakdown of importance sampling weights in its framework. This flaw induces substantial variance in the training process, leading to accumulated noise as response length increases. The clipping mechanism exacerbates this instability, frequently resulting in the collapse of the model.

To mitigate these fundamental challenges, we introduce \textbf{Group Sequence Policy Optimization (GSPO)}, a novel reinforcement learning approach tailored for large language model training. GSPO’s central contribution is a principled formulation of the importance ratio using sequence likelihood, drawing from theoretical underpinnings in the importance sampling literature \citep{click}. Moreover, GSPO incorporates normalized reward calculations based on the comparative advantages of multiple model responses per query, thereby fostering closer alignment between the reward system at the sequence level and the optimization process.

Our experimental results clearly indicate that GSPO substantially outperforms GRPO with respect to training stability, efficiency, and overall effectiveness. Notably, GSPO naturally addresses the stability issues typically encountered when training large Mixture-of-Experts (MoE) models in RL, thereby obviating the reliance on intricate stabilization methods and suggesting a pathway toward more streamlined RL infrastructures. The strengths of GSPO have played a pivotal role in achieving significant performance gains in recent Qwen3 models. We anticipate that GSPO will serve as a resilient and scalable algorithmic platform, facilitating ongoing progress in large-scale RL training for language models.

\section{Preliminaries}

\paragraph{Notation}
Throughout this work, we interpret an autoregressive language model with parameters $\theta$ as a policy $\pi_\theta$. The variable $x$ refers to a query, while the collection of all queries is represented by $\mathcal{D}$. For a given query $x$ and a generated response $y$, the conditional probability assigned by $\pi_\theta$ is expressed as $\pi_\theta (y | x)=\prod_{t=1}^{|y|} \pi_\theta (y_t | x, y_{<t} )$, where $|y|$ is the token count in $y$. A verifier $r$ assigns a reward $r(x, y)$ in the interval $[0, 1]$ to each query-response pair $(x, y)$.

\paragraph{Proximal Policy Optimization (PPO)} Using samples generated from the old policy $\pi_{\theta_\text{old}}$, PPO \citep{ppo} constrains the policy update within a proximal region of the old policy through the clipping mechanism. Specifically, PPO employs the following objective for policy optimization (we omit the KL regularization term hereinafter for brevity, as it is not the focus of this paper):

\begin{align}
\mathcal{J}_\text{PPO}(\theta) = \mathbb{E}_{ x \sim \mathcal{D},\, y \sim \pi_{\theta_\text{old}}( \cdot | x) }
\left[ \frac{1}{|y|} \sum_{t=1}^{|y|} 
\min \left( w_{t}(\theta) \widehat{A}_{t},  \, \mathrm{clip} \left( w_{t}(\theta), 1 - {\varepsilon}, 1 + {\varepsilon}\right) \widehat{A}_{t} \right)
\right],
\end{align}

In this context, the importance ratio for a token $y_t$ is given by $
w_{t}(\theta) = \frac{ \pi_{\theta} (y_{t} | x, y_{<t}) }{ \pi_{\theta_\text{old}} (y_{t} | x,y_{<t})} $, where the advantage $\widehat{A}_{t}$ of $y_t$ is approximated using a separate value model, and $\varepsilon$ serves as the threshold for clipping importance ratios.

A significant practical limitation of PPO stems from its strong dependence on the value model. Typically, the value model matches the policy model in size, resulting in substantial computational and memory overhead. Moreover, the overall success of the approach is critically dependent on the accuracy of the value estimate. Since obtaining a precise value model is already a difficult task, scaling this reliability to longer outputs and tasks of increased complexity is even more problematic.

\paragraph{Group Relative Policy Optimization (GRPO)} GRPO \citep{grpo} bypasses the need for the value model by computing the relative advantage of each response within a group of responses to the same query. Specifically, GRPO optimizes the following objective:

\begin{align}
\label{equ:grpo}
\mathcal{J}_\text{GRPO}(\theta) = \mathbb{E}_{ x \sim \mathcal{D},\, \{y_i\}_{i=1}^G \sim \pi_{\theta_\text{old}}( \cdot | x) }
\left[ \frac{1}{G} \sum_{i=1}^{G} \frac{1}{|y_i|} \sum_{t=1}^{|y_i|} 
\min \left( w_{i,t}(\theta) \widehat{A}_{i,t},  \, \mathrm{clip} \left( w_{i,t}(\theta), 1 - {\varepsilon}, 1 + {\varepsilon}\right) \widehat{A}_{i,t} \right)
\right],
\end{align}

where $G$ is the number of generated responses to each query $x$ (i.e., the group size), and the importance ratio $w_{i,t}(\theta)$ and advantage $\widehat{A}_{i,t}$ of token $y_{i,t}$ are:

\begin{align}
    w_{i,t}(\theta)=\frac{ \pi_{\theta} (y_{i,t} | x, y_{i,<t}) }{ \pi_{\theta_\text{old}} (y_{i,t} | x,y_{i,<t})},\quad\ 
    \widehat{A}_{i,t} = \widehat{A}_{i} = \frac{r(x, y_i) - \mathrm{mean} \left( \{ r(x, y_i) \}_{i=1}^G \right) }{ \mathrm{std} \left( \{ r(x, y_i) \}_{i=1}^G \right) },
\end{align}

respectively, where all the tokens in $y_i$ share the same advantage as $\widehat{A}_{i}$.

\section{Motivation}
\label{sec:motivation}

As model architectures become larger and more sparse—such as in Mixture-of-Experts frameworks—and as response outputs increase in length, reinforcement learning increasingly depends on processing large rollout batch sizes to fully leverage available hardware capabilities. To enhance the efficiency of sample usage, practitioners routinely divide these large rollout batches into several smaller mini-batches for the purpose of performing gradient updates. This mini-batch approach inherently results in an off-policy learning dynamic, since the sampled responses $y$ originate from a preceding policy, $\pi_{\theta_\text{old}}$, instead of the current optimization target, $\pi_\theta$. Consequently, this necessitates the use of the clipping strategy in algorithms like PPO and GRPO, as the clipping mechanism curbs the influence of samples that deviate substantially from the present policy during the calculation of gradients.

While mechanisms like clipping aim to manage this off-policy discrepancy, we identify a more fundamental issue in GRPO: \textit{its objective is ill-posed}. This problem becomes particularly acute when training large models on long-response tasks, leading to catastrophic model collapse. The ill-posed nature of the GRPO objective stems from a misapplication of importance sampling weights. The principle of importance sampling is to estimate the expectation of a function $f$ under a target distribution $\pi_\text{tar}$ by re-weighting samples drawn from a behavior distribution $\pi_\text{beh}$:

\begin{align}
\mathbb{E}_{z \sim \pi_\text{tar}} \left[ f(z) \right] = \mathbb{E}_{z \sim \pi_\text{beh}} \left[ \frac{ \pi_\text{tar} (z) }{ \pi_\text{beh} (z)} \, f(z) \right].
\end{align}

A key factor here is that effective adjustment for the discrepancy between distributions using the importance weight $\frac{ \pi_\text{tar} (z) }{ \pi_\text{beh} (z)}$ depends on taking an average over a substantial number of samples ($N \gg 1$) drawn from the behavior policy $\pi_\text{beh}$.

By contrast, GRPO employs an importance weight of $\frac{ \pi_{\theta} (y_{i,t} | x, y_{i,<t}) }{ \pi_{\theta_\text{old}} (y_{i,t} | x,y_{i,<t})}$ at every token position $t$. Because this weight corresponds to just a single sampled token $y_{i,t}$ from the subsequent-token distribution $\pi_{\theta_\text{old}}( \cdot | x, y_{i, <t})$, it cannot provide robust distributional correction. Instead, it introduces substantial variance into the gradient estimates during training—a problem that escalates with longer sequences and is further intensified by the presence of the clipping mechanism. In practice, this high-variance noise can precipitate a severe and sometimes permanent collapse of the model, from which recovery proves extremely difficult. Not even reverting to an earlier checkpoint, carefully adjusting hyperparameters (such as clipping ranges), extending the sequence length, or modifying RL targets can restore successful training after collapse.

These empirical findings highlight a foundational shortcoming in the architecture of GRPO. Specifically, the inadequacy of token-wise importance weighting reveals an essential tenet: \textbf{the granularity at which the optimization operates must align with the granularity at which rewards are allocated}. When rewards are provided at the sequence level, applying off-policy corrections at the token level is evidently suboptimal. This realization leads us to abandon the token-level approach in favor of directly applying importance weighting and optimization at the \textit{sequence level}.

\section{Algorithm}

\subsection{GSPO: Group Sequence Policy Optimization}

Although the use of token-level importance weighting $\frac{ \pi_{\theta} (y_{i,t} | x, y_{i,<t}) }{ \pi_{\theta_\text{old}} (y_{i,t} | x,y_{i,<t})}$ raises issues within GRPO, it is notable that, for language generation tasks, the \textit{sequence-level} importance weight $\frac{ \pi_{\theta} (y | x) }{ \pi_{\theta_\text{old}} (y | x)}$ possesses a well-defined theoretical interpretation: this metric quantifies the extent to which a generated response $y$—sampled according to $\pi_{\theta_\text{old}} (\cdot | x)$—diverges from the current model's distribution $\pi_{\theta} (\cdot | x)$. Consequently, this aligns well with rewards specified at the sequence level and functions as a useful measure for regulating the clipping process.

Based on this straightforward observation, we propose the \textbf{Group Sequence Policy Optimization (GSPO)} algorithm. GSPO employs the following sequence-level optimization objective:

\begin{align}
\mathcal{J}_\text{GSPO} (\theta) =
\mathbb{E}_{ x \sim \mathcal{D},\, \{y_i\}_{i=1}^G \sim \pi_{\theta_\text{old}}( \cdot | x) }
\left[ 
\frac{1}{G} \sum_{i=1}^{G}
\min \left( s_{i}(\theta)  \widehat{A}_{i},  \, \mathrm{clip} \left( s_{i}(\theta), 1 - {\varepsilon}, 1 + {\varepsilon} \right) \widehat{A}_{i} \right) 
\right],
\label{equ:gspo}
\end{align}

where we adopt the group-based advantage estimation:

\begin{align}
\widehat{A}_{i} = \frac{r(x, y_i) - \mathrm{mean} \left( \{ r(x, y_i) \}_{i=1}^G \right) }{ \mathrm{std} \left( \{ r(x, y_i) \}_{i=1}^G \right) },
\end{align}

and define the importance ratio $s_{i}(\theta)$ based on sequence likelihood \citep{click}:

\begin{align}
s_{i}(\theta) = \left( \frac{ \pi_{\theta} (y_i | x) }{ \pi_{\theta_\text{old}} (y_i | x)} \right)^{\frac{1}{|y_i|}}
=
\exp \left( \frac{1}{|y_i|} \sum_{t=1}^{|y_i|} \log \frac{ \pi_{\theta} (y_{i,t} | x, y_{i,<t}) }{ \pi_{\theta_\text{old}} (y_{i,t} | x,y_{i,<t})} \right).
\end{align}

Consequently, GSPO performs clipping at the sequence level rather than at the token level, thereby filtering out samples that are excessively ``off-policy'' from the gradient estimation process. This approach is consistent with both the reward computation, which is done per sequence, and the optimization procedure. It should also be emphasized that length normalization is incorporated into $s_{i}(\theta)$ to decrease variance and to ensure $s_{i}(\theta)$ remains within a consistent numerical interval. Failing to normalize by length can cause minor token likelihood changes to induce significant volatility in the overall sequence-level importance ratio, necessitating different clipping thresholds for sequences of various lengths. Moreover, we observe that the clipping bounds employed by GSPO typically have a different order of magnitude compared to those used in earlier methods such as GRPO, due to the fundamentally different ways the importance ratios are formulated.

\subsection{Gradient Analysis}
\label{subsec:gradient}

We can derive the gradient of the GSPO objective as follows (clipping is omitted for brevity):

\begin{align}
\nabla_{\theta} \mathcal{J}_\text{GSPO} (\theta)
=&\ 
\nabla_{\theta} \mathbb{E}_{ x \sim \mathcal{D},\, \{y_i\}_{i=1}^G \sim \pi_{\theta_\text{old}}( \cdot | x) }
\left[ 
\frac{1}{G} \sum_{i=1}^{G} s_{i}(\theta) \widehat{A}_{i}
\right] \\
= &\ 
\mathbb{E}_{ x \sim \mathcal{D},\, \{y_i\}_{i=1}^G \sim \pi_{\theta_\text{old}}( \cdot | x) }
\left[ 
\frac{1}{G} \sum_{i=1}^{G}
s_{i}(\theta) \widehat{A}_{i}
\cdot \nabla_{\theta} \log s_{i}(\theta)
\right] \\
=&\ 
\mathbb{E}_{ x \sim \mathcal{D},\, \{y_i\}_{i=1}^G \sim \pi_{\theta_\text{old}}( \cdot | x) }
\left[ 
\frac{1}{G} \sum_{i=1}^{G} 
\left( \frac{ \pi_{\theta} (y_i | x) }{ \pi_{\theta_\text{old}} (y_i | x)} \right)^{\frac{1}{|y_i|}} \widehat{A}_{i} 
\cdot \frac{1}{|y_i|} \sum_{t=1}^{|y_i|} \nabla_{\theta} \log \pi_{\theta} (y_{i,t} | x, y_{i,<t}) 
\right].
\label{equ:gspo_grad}
\end{align}

For comparison, the gradient of the GRPO objective is as follows (note that $\widehat{A}_{i,t} = \widehat{A}_{i}$):

\begin{align}
\nabla_{\theta} \mathcal{J}_\text{GRPO} (\theta) 
=&\ 
\nabla_{\theta} \mathbb{E}_{ x \sim \mathcal{D},\, \{y_i\}_{i=1}^G \sim \pi_{\theta_\text{old}}( \cdot | x) }
\left[ \frac{1}{G} \sum_{i=1}^{G} \frac{1}{|y_i|} \sum_{t=1}^{|y_i|} 
w_{i,t}(\theta) \widehat{A}_{i,t}
\right] \\
=&\
\mathbb{E}_{ x \sim \mathcal{D},\, \{y_i\}_{i=1}^G \sim \pi_{\theta_\text{old}}( \cdot | x) }
\left[ \frac{1}{G} \sum_{i=1}^{G} \widehat{A}_{i} 
\cdot \frac{1}{|y_i|} \sum_{t=1}^{|y_i|} 
\frac{ \pi_{\theta} (y_{i,t} | x, y_{i,<t}) }{ \pi_{\theta_\text{old}} (y_{i,t} | x,y_{i,<t})} 
\nabla_{\theta} \log \pi_{\theta} (y_{i,t} | x, y_{i,<t})  
\right].
\end{align}

Consequently, the key difference between GSPO and GRPO concerns \textit{the manner in which they apply weights to the gradients of the tokens’ log likelihoods}. GRPO employs importance weighting for tokens, assigning each a weight of $\frac{ \pi_{\theta} (y_{i,t} | x, y_{i,<t}) }{ \pi_{\theta_\text{old}} (y_{i,t} | x,y_{i,<t})}$. These token-specific weights can vary significantly—taking values in $(0, 1+\varepsilon]$ for $\widehat{A}_{i} >0$ or $[1-\varepsilon, +\infty)$ when $\widehat{A}_{i} < 0$—and are not trivial. As training advances, the cumulative effect of these disparities can introduce instability and undesirable dynamics. By contrast, GSPO uniformly weights all tokens within a response, thereby removing the potential for the instability present in GRPO.

\subsection{GSPO-token: A Token-level Objective Variant}

In scenarios like multi-turn RL, we may desire a finer-grained advantage adjustment than the sequence level. To this end, we introduce a token-level objective variant of GSPO, namely \textbf{GSPO-token}, to allow token-wise advantage customization:

\begin{align}
\mathcal{J}_\text{GSPO-token}(\theta)
=
\mathbb{E}_{ x \sim \mathcal{D},\, \{y_i\}_{i=1}^G \sim \pi_{\theta_\text{old}}( \cdot | x) }
\left[ 
\frac{1}{G} \sum_{i=1}^{G} \frac{1}{|y_i|} \sum_{t=1}^{|y_i|} 
\min \left( s_{i,t}(\theta) \widehat{A}_{i,t},  \, \mathrm{clip} \left( s_{i,t}(\theta), 1 - {\varepsilon}, 1 + {\varepsilon} \right) \widehat{A}_{i,t} \right) 
\right],
\label{equ:gspo-token}
\end{align}

where

\begin{align}
s_{i,t}(\theta) 
= \mathrm{sg} \left[ s_{i}(\theta) \right]  \cdot \frac{ \pi_{\theta} (y_{i,t} | x, y_{i,<t}) }{ \mathrm{sg} \left[ \pi_{\theta} (y_{i,t} | x, y_{i,<t}) \right] },
\end{align}

and $\mathrm{sg}[\cdot]$ denotes only taking the numerical value but stopping the gradient, corresponding to the \texttt{detach} operation in PyTorch. The gradient of GSPO-token can be derived as:

\begin{align}
\nabla_{\theta} \mathcal{J}_\text{GSPO-token}(\theta)
=&\ 
\nabla_{\theta} \mathbb{E}_{ x \sim \mathcal{D},\, \{y_i\}_{i=1}^G \sim \pi_{\theta_\text{old}}( \cdot | x) }
\left[ 
\frac{1}{G} \sum_{i=1}^{G}
\frac{1}{|y_i|} \sum_{t=1}^{|y_i|} 
s_{i,t}(\theta) \widehat{A}_{i,t}
\right] \\
=&\ 
\mathbb{E}_{ x \sim \mathcal{D},\, \{y_i\}_{i=1}^G \sim \pi_{\theta_\text{old}}( \cdot | x) }
\left[ 
\frac{1}{G} \sum_{i=1}^{G} s_i (\theta) \cdot \frac{1}{|y_i|} \sum_{t=1}^{|y_i|} 
\widehat{A}_{i,t} \frac{ \nabla_{\theta} \pi_{\theta} (y_{i,t} | x, y_{i,<t}) }{ \pi_{\theta} (y_{i,t} | x, y_{i,<t}) }
\right] \\
=&\ 
\mathbb{E}_{ x \sim \mathcal{D},\, \{y_i\}_{i=1}^G \sim \pi_{\theta_\text{old}}( \cdot | x) }
\left[ 
\frac{1}{G} \sum_{i=1}^{G}  \left( \frac{ \pi_{\theta} (y_i | x) }{ \pi_{\theta_\text{old}} (y_i | x)} \right)^{\frac{1}{|y_i|}}
\cdot \frac{1}{|y_i|} \sum_{t=1}^{|y_i|}
\widehat{A}_{i,t} \nabla_{\theta} \log \pi_{\theta} (y_{i,t} | x, y_{i,<t})  
\right].
\label{equ:gspo-token_grad}
\end{align}

Observe that $\frac{ \pi_{\theta} (y_{i,t} | x, y_{i,<t}) }{ \mathrm{sg} \left[ \pi_{\theta} (y_{i,t} | x, y_{i,<t}) \right] }$ evaluates to 1, which implies that $s_{i,t}(\theta)$ is, in effect, equivalent to $s_{i}(\theta)$ in terms of numeric value.
When analyzing Equation~\eqref{equ:gspo} alongside~\eqref{equ:gspo-token}, as well as Equation~\eqref{equ:gspo_grad} with~\eqref{equ:gspo-token_grad}, it becomes clear that GSPO-token and GSPO yield identical numerical outcomes for the optimization objective, clipping condition, and theoretical gradient provided that all token-level advantages within a response $y_i$ are set uniformly (that is, $\widehat{A}_{i,t} = \widehat{A}_{i}$ for all $t$). Nevertheless, GSPO-token offers increased adaptability by allowing the assignment of distinct advantage values for each token.

\section{Experiments and Discussion}

\subsection{Empirical Results}

We conduct experiments utilizing a cold-start model initialized from Qwen3-30B-A3B-Base, presenting both the progression of training rewards and the resulting model performance on the AIME'24 dataset (computed as average Pass@1 across 32 samples), LiveCodeBench (202410-202502, average Pass@1 across 8 samples), and CodeForces (evaluated by Elo Rating). Throughout the RL optimization process, each rollout batch is subdivided into four mini-batches to perform gradient steps. For GSPO, the left and right clipping bounds in Equation~\eqref{equ:gspo} are chosen to be 3e-4 and 4e-4, respectively. For the baseline, we utilize GRPO, assigning the clipping ranges in Equation~\eqref{equ:grpo} as 0.2 on the left and 0.27 on the right; these hyperparameters have been thoroughly tuned for equivalency in comparison. It is important to highlight that GRPO requires application of the Routing Replay training method to achieve standard convergence behavior with MoE RL—a requirement that will be further elaborated in \S~\ref{subsec:routing_replay}—whereas \textit{GSPO eliminates the reliance on this mechanism}.

Figure~\ref{fig:results} illustrates that training with GSPO maintains stability across the entire process. Notably, \textbf{GSPO consistently enhances performance as training compute is scaled up, the query set is periodically refreshed, and the generation length is increased}. In addition, GSPO outperforms GRPO in terms of training efficiency, reaching higher training accuracy and benchmark scores given identical training compute and query consumption. Finally, our successful application of GSPO to the RL training of state-of-the-art Qwen3 models provides compelling evidence of GSPO’s effectiveness in leveraging RL scaling for large language models.

\subsection{Curious Observation on Clipping Fractions}

An important differentiator between GSPO and GRPO is that GSPO clips responses at the sequence level, rather than at the token level. As depicted in Figure~\ref{fig:clipping}, this methodology results in a difference of approximately two orders of magnitude in the proportions of tokens that are clipped under GSPO versus GRPO (a disparity that persists even when the clipping ranges are varied). Notably, although GSPO clips many more tokens overall—resulting in fewer data points available for training or for the purpose of estimating gradients—it consistently demonstrates greater training efficiency compared to GRPO. This somewhat unexpected outcome—wherein removing a substantially larger share of tokens still yields improved training efficiency—suggests that GRPO’s token-level gradient computations suffer from higher intrinsic noise and offer limited effectiveness in exploiting training samples. By contrast, GSPO’s reliance on sequence-level information provides a stronger and more dependable learning signal.

\subsection{Benefit of GSPO for MoE Training}
\label{subsec:routing_replay}

\paragraph{Background}
In contrast to RL training with dense architectures, MoE models present distinctive stability issues due to their sparsely activated pathways. Notably, during experiments utilizing the GRPO algorithm, we observed that MoE models exhibit substantial \textit{expert-activation volatility}, which can disrupt the convergence of RL optimization. Specifically, following gradient updates, the set of experts selected for generating identical responses may alter considerably. Illustratively, with the 48-layer Qwen3-30B-A3B-Base configuration, analysis reveals that, after each RL gradient step and considering the same rollout instance, nearly 10\% of the experts chosen according to the updated policy $\pi_{\theta}$ differ from those picked by the previous policy $\pi_{\theta_\text{old}}$. This behavior—especially pronounced in larger and deeper MoE models—leads to significant instability in the token-level importance ratios $w_{i,t}(\theta) = \frac{ \pi_{\theta} (y_{i,t} | x, y_{i,<t}) }{ \pi_{\theta_\text{old}} (y_{i,t} | x,y_{i,<t})}$, causing these ratios to become highly erratic and, as detailed in \S~\ref{sec:motivation} and~\ref{subsec:gradient}, ultimately unreliable. As a result, this instability impedes effective RL training convergence.

\paragraph{Our Previous Approach} In our earlier work, we addressed this issue by utilizing the \textbf{Routing Replay} training method. This involves storing the sets of experts activated by $\pi_{\theta_\text{old}}$ and later reenacting these exact routing decisions when evaluating $\pi_{\theta}$ to compute the importance ratios $w_{i,t}(\theta) = \frac{ \pi_{\theta} (y_{i,t} | x, y_{i,<t}) }{ \pi_{\theta_\text{old}} (y_{i,t} | x,y_{i,<t})}$. Through this mechanism, for any token $y_{i,t}$, both $\pi_{\theta} (y_{i,t} | x, y_{i,<t})$ and $\pi_{\theta_\text{old}} (y_{i,t} | x, y_{i,<t})$ utilize an identical configuration of expert activations, thereby preserving the reliability of the computed token-level importance ratios. This approach ensures that gradient-based learning consistently updates the same underlying subnetwork. As shown in Figure~\ref{fig:routing_replay}, Routing Replay plays a pivotal role in maintaining stable and effective convergence during the GRPO optimization of MoE models.

\paragraph{Benefit of GSPO} While Routing Replay allows GRPO to successfully train MoE models by ensuring convergence, its approach of recycling previous routing patterns leads to increased memory and communication costs and may artificially cap the model's usable capacity. On the other hand, as depicted in Figure~\ref{fig:results}, GSPO removes the reliance on Routing Replay, permitting standard computation of importance ratios $s_i(\theta)$, and consequently achieves normal convergence and stable optimization. The central idea behind GSPO is its exclusive focus on sequence likelihoods (i.e., $\pi_{\theta} (y_i | x)$), rather than on per-token likelihoods (i.e., $\pi_{\theta} (y_{i,t} | x, y_{i,<t})$), thereby rendering it insensitive to fluctuations at the token level. Given that the MoE model inherently preserves its language modeling function, the sequence likelihood values remain relatively stable. To conclude, GSPO addresses the issue of expert-activation instability at its root, thereby making auxiliary techniques like Routing Replay unnecessary. As a result, training becomes more streamlined and reliable, and the model can fully exploit its architecture without imposed limitations.

\subsection{Benefit of GSPO for RL Infrastructure}

Owing to differences in numerical precision between training engines (such as Megatron) and inference engines (such as SGLang and vLLM), it is common practice to recalculate the likelihoods of sampled outputs using the training engine when evaluating under the previous policy $\pi_{\theta_\text{old}}$. In contrast, GSPO relies on sequence-level likelihoods for its optimization process rather than token-level likelihoods, and the former generally exhibits greater robustness to precision mismatches. As a result, GSPO enables the direct use of likelihood values produced by the inference engine during optimization, eliminating the necessity for additional computations using the training engine. This property proves advantageous in contexts such as partial rollout, multi-turn RL, and architectures where training and inference are decoupled.

\section{Conclusion}

We introduce Group Sequence Policy Optimization (GSPO), a novel reinforcement learning algorithm designed for large language model training. Grounded in the fundamental concepts of importance sampling, GSPO employs sequence likelihood to calculate importance ratios and applies sequence-level procedures for clipping, reward assignment, and policy optimization. Compared to GRPO, GSPO delivers markedly better stability, training efficiency, and model performance, proving especially effective in large-scale RL scenarios involving MoE models. This serves as a crucial enabler for the significant advancements observed in recent Qwen3 models. Leveraging GSPO's scalability, we intend to further expand the reach of RL, anticipating substantial progress in artificial intelligence.

\end{document}